\documentclass[12pt,letterpaper]{article}

\usepackage[utf8]{inputenc}
\usepackage[T1]{fontenc}
\usepackage[spanish,es-noquoting,es-tabla]{babel}
\usepackage[margin=2cm]{geometry}
\usepackage{carlito}
\renewcommand{\familydefault}{\sfdefault}
\usepackage{setspace}
\usepackage{tabularx}
\usepackage{array}
\usepackage[table]{xcolor}
\usepackage{booktabs}
\usepackage{parskip}
\usepackage{enumitem}
\usepackage[hidelinks]{hyperref}

\definecolor{evmblue}{HTML}{4F9AC8}
\setlength{\parindent}{0pt}
\setlength{\parskip}{0.3em}
\setlength{\tabcolsep}{3pt}
\renewcommand{\arraystretch}{1.2}
\newcolumntype{C}[1]{>{\centering\arraybackslash}p{#1}}
\setlist{nosep,leftmargin=1.6em}

\begin{document}
\onehalfspacing

\begin{center}
  {\large\bfseries Método de Gestión del Valor Ganado (EVM)}\\[0.6em]
{\large\bfseries Carlos Sebastián Madrigal Rodríguez}\\[0.8em]
\end{center}

Se parte de un proyecto de desarrollo de software
cuya duración estimada es de seis meses y cuyo costo total estimado es de
\$115\,000. Los valores planificados y acumulados para cada mes del proyecto son
los siguientes:

\begin{table}[htbp]
\centering
\small
\begin{tabularx}{\textwidth}{|>{\raggedright\arraybackslash}X|*{7}{C{1.7cm}|}}
\hline
\rowcolor{evmblue}
\color{white}\textbf{Meses} &
\color{white}\textbf{1} &
\color{white}\textbf{2} &
\color{white}\textbf{3} &
\color{white}\textbf{4} &
\color{white}\textbf{5} &
\color{white}\textbf{6} &
\color{white}\textbf{Total} \\
\hline
Línea base de costo del proyecto (PV) &
\$23\,000 & \$15\,000 & \$16\,000 & \$19\,000 & \$20\,000 & \$22\,000 & \$115\,000 \\
\hline
\rowcolor{evmblue!12}
Valores acumulados de PV &
\$23\,000 & \$38\,000 & \$54\,000 & \$73\,000 & \$93\,000 & \$115\,000 & \$115\,000 \\
\hline
\end{tabularx}
\end{table}

\section*{Parte A. Control del proyecto en el mes 3}

Teniendo como referencia esta línea base, supongamos que el proyecto se
encuentra en el mes 3 de ejecución. Contamos con mediciones reales del trabajo
realizado y del costo real hasta ese mes, cuyos valores se recogen en la
siguiente tabla:

\begin{table}[htbp]
\centering
\begin{tabularx}{\textwidth}{|>{\raggedright\arraybackslash}X|*{7}{C{1.7cm}|}}
\hline
\rowcolor{evmblue}
\color{white}\textbf{Meses} &
\color{white}\textbf{1} &
\color{white}\textbf{2} &
\color{white}\textbf{3} &
\color{white}\textbf{4} &
\color{white}\textbf{5} &
\color{white}\textbf{6} &
\color{white}\textbf{Total} \\
\hline
Costo real (AC) &
\$21\,000 & \$18\,000 & \$19\,000 & -- & -- & -- & \$58\,000 \\
\hline
\rowcolor{evmblue!12}
Valores acumulados de AC &
\$21\,000 & \$39\,000 & \$58\,000 & -- & -- & -- & \$58\,000 \\
\hline
Avance real (trabajo realizado) &
8\,\% & 15\,\% & 14\,\% & -- & -- & -- & 37\,\% \\
\hline
\end{tabularx}
\end{table}

\subsection*{Indicadores de valor ganado}

Para el análisis se considera un presupuesto al término (BAC) de \$115\,000.
Al cierre del mes 3, el valor planificado acumulado (PV) es \$54\,000 y el
costo real acumulado (AC) es \$58\,000. El avance real acumulado es la suma
de 8\%, 15\% y 14\%, es decir, 37\%, por lo que el valor ganado (EV) es:

\[
EV = BAC \times 0{,}37 = \$115\,000 \times 0{,}37 = \$42\,550.
\]

\begin{itemize}
  \item \textbf{Índice de desempeño del costo (CPI):}
  $EV/AC = 42\,550/58\,000 \approx 0{,}73$.
  \item \textbf{Índice de desempeño del cronograma (SPI):}
  $EV/PV = 42\,550/54\,000 \approx 0{,}79$.
  \item \textbf{Variación del costo (CV):}
  $EV-AC = 42\,550-58\,000 = -\$15\,450$.
  \item \textbf{Variación del cronograma (SV):}
  $EV-PV = 42\,550-54\,000 = -\$11\,450$.
\end{itemize}

\subsection*{Respuestas}

\begin{enumerate}
  \item \textbf{¿El proyecto va por encima o por debajo del presupuesto?}
  Va por encima del presupuesto para el trabajo realizado. El CPI es menor que
  1 y la variación del costo es negativa: se han gastado \$58\,000 en trabajo
  cuyo valor presupuestado es \$42\,550.

  \item \textbf{¿El proyecto va retrasado o adelantado?}
  Va retrasado. El SPI es menor que 1 y el trabajo completado tiene un valor
  \$11\,450 inferior al que estaba planificado para el mes 3.

  \item \textbf{¿Cuál será el costo estimado al finalizar el proyecto a este
  ritmo?}
  Si se mantiene el desempeño actual del costo, la estimación al término (EAC)
  es:
  \[
  EAC = BAC/CPI = \$115\,000/(42\,550/58\,000) \approx \$156\,757.
  \]
  Esto representa un sobrecosto estimado de aproximadamente \$41\,757 frente
  al presupuesto original.

  \item \textbf{¿Cuándo se estima finalizar el proyecto a este ritmo?}
  Usando el SPI actual como proyección lineal del desempeño del cronograma, la
  duración estimada es:
  \[
  6/SPI = 6/0{,}7880 \approx 7{,}61.
  \]
  Por tanto, se estima una duración de 7,61 meses y un retraso de
  aproximadamente 1,61 meses respecto de
  los seis meses previstos. Es una proyección aproximada que supone que el SPI
  se mantiene durante el resto del proyecto.

  \item \textbf{¿Qué se podría hacer para controlar el proyecto?}
  \begin{itemize}
    \item \textbf{Recuperar el desempeño del costo:} revisar las variaciones
    por paquete de trabajo para identificar sus causas; reforzar el control de
    cambios para evitar trabajo no aprobado y renegociar compras o asignaciones
    que no sean críticas. Actualizar el EAC después de aplicar las medidas.
    \item \textbf{Reducir el retraso:} revisar la ruta crítica y los bloqueos,
    reasignar personal con las habilidades necesarias a las actividades
    críticas y evaluar la ejecución en paralelo de tareas que no dependan
    entre sí. Antes de acelerar, valorar el costo adicional y los riesgos para
    la calidad.
    \item \textbf{Intensificar el seguimiento:} actualizar PV, EV y AC con
    mayor frecuencia, revisar CPI y SPI en cada corte y asignar responsables y
    fechas a las acciones correctivas.
  \end{itemize}
\end{enumerate}

\section*{Parte B: Mes 5}

En el mes 5 se cuenta con las siguientes mediciones reales:

\begin{table}[htbp]
\centering
\footnotesize
\begin{tabularx}{0.95\textwidth}{|>{\raggedright\arraybackslash}X|*{7}{C{1.5cm}|}}
\hline
\rowcolor{evmblue} \color{white}\textbf{Meses} &
\color{white}\textbf{1} &
\color{white}\textbf{2} &
\color{white}\textbf{3} &
\color{white}\textbf{4} &
\color{white}\textbf{5} &
\color{white}\textbf{6} &
\color{white}\textbf{Total} \\
\hline
Costo real (AC) & \$21k & \$18k & \$19k & \$20k & \$21k & -- & \$99k \\
\hline
Acum.\ AC & \$21k & \$39k & \$58k & \$78k & \$99k & -- & \$99k \\
\hline
Avance & 8\,\% & 15\,\% & 14\,\% & 21\,\% & 22\,\% & -- & 80\,\% \\
\hline
\end{tabularx}
\end{table}

Al cierre del mes 5, PV acumulado es \$93\,000 (línea base) y AC acumulado es
\$99\,000. El avance real acumulado es 8\,\% + 15\,\% + 14\,\% + 21\,\% +
22\,\% = 80\,\%, por lo que:

\[
EV = BAC \times 0{,}80 = \$115\,000 \times 0{,}80 = \$92\,000.
\]

\textbf{Indicadores:}
\begin{itemize}
  \item $CPI = EV/AC = 92\,000/99\,000 \approx 0{,}93$.
  \item $SPI = EV/PV = 92\,000/93\,000 \approx 0{,}99$.
  \item $CV = EV - AC = 92\,000 - 99\,000 = -\$7\,000$.
  \item $SV = EV - PV = 92\,000 - 93\,000 = -\$1\,000$.
  \item $EAC = BAC/CPI = 115\,000/0{,}93 \approx \$123\,656$.
  \item Duración: $6/SPI = 6/0{,}99 \approx 6{,}06$ meses.
\end{itemize}

\subsection*{Respuestas}

\begin{enumerate}
  \item \textbf{¿El proyecto va por encima o por debajo del presupuesto?}
  Va por encima: $CPI=0{,}93<1$ y $CV=-\$7\,000<0$. Se ha gastado más de lo
  presupuestado para el trabajo completado, aunque la mejora respecto al mes 3
  es significativa.

  \item \textbf{¿El proyecto va retrasado o adelantado?}
  Va ligeramente retrasado: $SPI=0{,}99<1$ y $SV=-\$1\,000<0$. El valor del
  trabajo completado es inferior al planificado, pero el retraso es mínimo
  comparado con el mes 3.

  \item \textbf{¿Cuál será la estimación de costo al finalizar?}
  $EAC\approx\$123{,}656$, un sobrecosto estimado de \$8{,}656. Mejora
  sustancial frente a los \$156{,}757 proyectados en el mes 3.

  \item \textbf{¿Se finalizará el proyecto con retraso o adelanto?}
  Con retraso: la duración estimada es 6,06 meses, aproximadamente 0,06 meses
  (aproximadamente 2 días) después de lo planificado. Es un retraso mínimo
  comparado con los 1,61 meses del mes 3.

  \item \textbf{¿Las medidas tomadas en el mes 3 tuvieron efecto positivo o
  negativo?}
  Tuvieron efecto \textbf{positivo}. Los indicadores mejoraron notablemente:
  CPI aumentó de 0,73 a 0,93 (28\,\%), SPI de 0,79 a 0,99 (25\,\%), CV de
  -\$15{,}45k a -\$7k (55\,\% de mejora) y SV de -\$11{,}45k a -\$1k (91\,\%
  de mejora). Aunque no es prueba definitiva de causalidad, la evolución
  consistente y sustancial en dos meses sugiere que las acciones correctivas
  implementadas fueron efectivas en reducir desviaciones de costo y cronograma.
\end{enumerate}

\end{document}